%% kapa: exam=1 round=25 year=2014 period=1 form=A subject=civil_law subject_ko=민법 seq=8 number=8
%% source: Q-Net 감정평가사 기출문제 (cst00309) artlSeq=5101207 file=extracted/1008/5101207/2014년 제25회 감정평가사 1차시험 기출문제/1교시(A형).hwp
%% conversion: kordoc 4.24.51 -> tools/gen1.py (auto); review=proofread-sonnet; flags=none
\begin{question}{8}
물건에 관한 설명으로 옳은 것은? (다툼이 있으면 판례에 의함)
\begin{choices}
  \item 부동산은 언제나 종물이 될 수 없다.
  \item 주물 자체의 효용과 직접적인 관계가 없는 물건도 종물이 될 수 있다.
  \item 물건의 용법에 의하여 수취하는 산출물은 법정과실에 해당한다.
  \item 천연과실은 수취할 권리의 존속기간 일수의 비율로 취득함이 원칙이다.
  \item 건물의 개수는 공부상의 등록뿐만 아니라 거래관념에 따라 제반사정을 참작하여 결정한다.
\end{choices}
\end{question}
\answer{5}
